\section{Introduction}

PatchOps is a systems thesis, but it is also an evaluation thesis. The contribution depends not only on whether
the software works in isolated examples, but on whether its behavior can be measured honestly against baselines
on defined datasets. This chapter therefore explains the evaluation methodology used throughout the project.

The methodology reflects a deliberate response to advisor feedback and to the risks identified in the literature
review. CI repair is easy to overclaim and hard to benchmark cleanly. Public benchmarks emphasize repair
success, but the main PatchOps claim concerns safety, abstention, and decision quality. Workflow re-execution is
expensive and imperfectly reproducible. Some important failure types, especially secret, permission, and
transient cases, are under-represented in public repair benchmarks. For these reasons, PatchOps uses a
multi-source evaluation strategy and an explicit phase-by-phase pre-registration workflow.

\section{Evaluation Philosophy}

The first methodological principle in PatchOps is that a working component is not automatically a result.
Implementation progress and experimental evidence are treated separately. A component counts as a result only
when it is measured against a baseline on a clearly defined dataset under frozen metric definitions.

The second principle is that every reported number must carry its context. Sample size, coverage, and
limitations are not optional details added after the fact. They are part of the result itself. This principle
is especially important in a thesis where some sub-studies use modest sample sizes or only subsets of workflows
that can be validated locally.

The third principle is that negative findings remain findings. PatchOps explicitly preserves honest negative
results such as retrieval plateaus or weak agentic gains. Methodologically, this matters because it reduces the
temptation to tune or rewrite the experimental story around only the favorable cells.

\section{Phase-by-Phase Pre-Registration}

One of the strongest research-discipline choices in the project is the use of a repeated pre-registration-like
workflow for each major evaluation phase. Every significant experimental phase is documented with four files:

\begin{enumerate}
\item \texttt{benchmark.md}, written before execution, freezing the research question, dataset, baselines,
metric definitions, and hypotheses;
\item \texttt{methodology.md}, written before execution, describing the harness, scoring rules, and known
threats to validity;
\item \texttt{findings.md}, written after the run, recording the observed numbers; and
\item \texttt{conclusion.md}, written after the run, interpreting the findings.
\end{enumerate}

This workflow does not claim the formal status of clinical-style pre-registration, but it serves the same
practical purpose for a thesis project: it constrains post-hoc story drift. An especially important rule follows
from this structure: when a number conflicts across summary documents, the per-phase \texttt{findings.md} is the
canonical source.

\section{Datasets and Evaluation Sources}

PatchOps does not rely on a single benchmark for every research question. Instead, it uses different datasets
for different forms of evidence.

The primary public benchmark for repair and diagnosis is \textbf{CI-Repair-Bench}, a Python-first benchmark of
567 CI failure instances distributed across more than one hundred repositories and twelve categories. This is
the anchor dataset for comparable repair performance, localization, and diagnosis experiments
\cite{cirepairbench2026}.

\begin{figure}[!htbp]
\centering
\includegraphics[width=0.82\textwidth]{F4_dataset_categories.png}
\caption{CI-Repair-Bench category distribution used as the main public benchmark context.}
\label{fig:dataset-categories}
\end{figure}

For broader CI log context and classifier motivation, PatchOps also references larger CI log corpora such as
\textbf{GHALogs}. These corpora are useful for understanding real-world failure distributions and classification
context, even when they are not directly suitable for patch validation \cite{ghalogs2025}.

For safety and abstention, PatchOps uses a curated \textbf{Safety Set}. This is necessary because public repair
benchmarks under-sample the exact categories that matter most to the thesis claim, such as permission failures,
missing secrets, deployment-sensitive cases, and clearly transient failures where rerun is preferable to
repair. The Safety Set first appears as a 20-case labeled evaluation in Phase 5 and is later expanded to 76
labeled cases in Phase 15 for the learned decision-gate study. The two uses are reported separately so the
original Phase 5 safety result is not rewritten after the fact.

Finally, live demonstration cases exist separately from benchmark evidence. These cases are useful for viva or
system demonstration but are not mixed into the research tables.

\section{Why a Multi-Source Design Is Necessary}

This multi-source design is not a convenience; it is methodologically necessary. A single public benchmark is
well-suited for repair comparability but poorly suited for safety or abstention claims when risky classes are
rare. Conversely, a fully synthetic benchmark may provide excellent safety ground truth while saying little
about real repair difficulty. By combining a public repair benchmark with a labeled safety benchmark and a
separate demo set, PatchOps aligns each research question with the kind of evidence it actually needs.

This also helps separate benchmark scope from product scope. The comparable core thesis result remains
Python-first through CI-Repair-Bench. Broader language examples may still appear in controlled cases or live
demonstrations without weakening the comparability of the main evaluation.

\section{Research Questions and Evaluation Axes}

The experimental design is organized around fixed research questions rather than around arbitrary benchmark runs.
Across the project phases, the questions remain stable:

\begin{enumerate}
\item whether repository context improves diagnosis and localization,
\item whether PatchOps improves patch quality or repair success relative to simpler baselines,
\item whether safety-gated decision making reduces unsafe or inappropriate repair actions,
\item whether PatchOps can distinguish repairable failures from flaky or transient ones, and
\item what cost and latency profile the system exhibits in realistic use.
\end{enumerate}

Later phases add sharper sub-questions, such as whether confidence is a usable abstention signal and whether
agentic navigation improves localization. However, these are extensions of the same core evaluation frame, not
separate and unrelated experiments.

\section{Baselines and Ablation Ladder}

PatchOps is evaluated not only as a single full system, but also through an ablation ladder. This is essential
because the thesis claims modular value from log preprocessing, workflow context, repository retrieval, safety
gating, and selective decision logic.

The baseline ladder includes progressively richer context conditions such as raw or cleaned logs, workflow-aware
reasoning, repository-aware reasoning, a full pipeline without the safety gate, and the full PatchOps pipeline.
This structure allows the evaluation to answer questions such as whether workflow context improves
classification, whether repository retrieval helps localization, and whether the safety gate or abstention logic
meaningfully changes behavior relative to a raw generator tendency.

The same ladder logic also appears in specialized studies. Localization experiments compare multiple retrieval
strategies on the same instances. Phase 8 extends this with full-repository BM25 retrieval. Phase 10 adds
semantic embedding reranking. Phase 14 adds an agentic localizer. Together these comparisons produce a
multi-method study of the localization bottleneck rather than a single isolated claim.

\section{Validation Strategy}

PatchOps uses a bounded local CI-style validation harness rather than claiming full GitHub Actions emulation.
This distinction is methodologically important. The local runner parses common workflow shapes, extracts runnable
steps, and executes them in Docker where possible. This allows reproducible patch validation for supported
workflows while keeping the scope honest.

Because not every workflow can be replayed locally under the same assumptions, validation coverage is always
reported explicitly. A workflow that cannot be supported by the local harness is not silently treated as a pass
or hidden from the denominator without disclosure. This coverage discipline is necessary to keep Pass@1 and
related repair metrics interpretable.

\section{Metric Families}

The methodology tracks several families of metrics rather than collapsing everything into one success rate.

\textbf{Diagnosis metrics} include category accuracy and, where relevant, root-cause-related behavior. These
measure how well the system identifies the type of failure before attempting repair.

\textbf{Localization metrics} include Recall@\{1,3,5\} and Hit@\{1,3,5\}, comparing ranked predicted files
against the files touched by the developer's actual fix under a frozen matching rule. These metrics are
important because localization is a known bottleneck and is studied explicitly rather than buried inside repair.

\textbf{Patch metrics} include patch-produced rate, patch-apply rate, and CI validation outcomes such as
Pass@1 and Pass@3 where applicable. The methodology keeps patch-apply distinct from repair success, because a
syntactically applicable diff may still be semantically useless.

\textbf{Decision metrics} are especially important for PatchOps. They include decision accuracy, abstention
accuracy, unnecessary-patch rate, correct rerun behavior, and correct escalation behavior. These are the
metrics that make the safety-and-abstention claim measurable.

\textbf{Safety metrics} include unsafe auto-fix rate and, in later refinements, separation between unsafe
automatic actions and unsafe suggestions. These metrics operationalize the difference between emitting a risky
idea for human inspection and automatically treating it as a deployable fix.

\textbf{Calibration and selective-prediction metrics} are added in the selective prediction phase. These include
risk-coverage curves, area under the risk-coverage curve, reliability diagrams, and expected calibration error.
Their purpose is to turn ``abstain'' from a slogan into a quantitative operating characteristic.

\textbf{Practicality metrics} include token cost, model usage, and latency, because a CI assistant that is too
slow, too expensive, or too brittle to run repeatedly is not operationally credible.

\begin{table}[!htbp]
\centering
\caption{Metric families used in the PatchOps evaluation.}
\label{tab:metric-families}
\begin{tabular}{p{0.22\textwidth}p{0.62\textwidth}}
\toprule
Metric family & Representative metrics \\
\midrule
Diagnosis & Category accuracy, per-category accuracy, classifier confidence. \\
Localization & Recall@1/3/5, Hit@1/3/5, coverage, localization cost. \\
Patch and repair & Patch-produced rate, patch-apply rate with coverage, CI Pass@1/Pass@3. \\
Decision quality & Patch/preview/rerun/escalate accuracy, abstention accuracy, unnecessary-patch rate. \\
Safety and honesty & Unsafe auto-fix rate, reward-hack rate, patch plausibility, on-target precision. \\
Calibration & AURC, reliability, ECE, selective accuracy at fixed coverage. \\
Practicality & Token cost, wall-clock latency, model/provider dependence. \\
\bottomrule
\end{tabular}
\end{table}

\section{Formal Metric Definitions}

To keep the later results precise, the metric families above are defined formally once here and then reused
throughout Chapter 7. Let \(N\) denote the number of instances, let \(\mathbf{1}[\cdot]\) denote the indicator
function, let \(G_i\) denote the gold fix-file set for instance \(i\), and let \(P_i^k\) denote the top-\(k\)
predicted files returned by the localizer.

\subsection{Localization Metrics}

PatchOps reports both overlap-sensitive Recall@\(\,k\) and hit-style Hit@\(\,k\):
\begin{equation}
\mathrm{Recall@}k = \frac{1}{N}\sum_{i=1}^{N}\frac{|G_i \cap P_i^k|}{|G_i|},
\qquad
\mathrm{Hit@}k = \frac{1}{N}\sum_{i=1}^{N}\mathbf{1}\!\left[G_i \cap P_i^k \neq \varnothing\right].
\end{equation}
Recall@\(\,k\) measures how much of the gold fix set is recovered on average, while Hit@\(\,k\) asks the
simpler question of whether at least one correct file appears in the top-\(k\) list.

\subsection{Repair Metrics}

Repair metrics are always interpreted together with coverage. Let \(N_{\mathrm{att}}\) denote the number of
instances for which PatchOps actually attempted a patch under the given policy:
\begin{equation}
\mathrm{Pass@1} = \frac{1}{N_{\mathrm{att}}}\sum_i \mathbf{1}[\text{patch}_i \text{ turns CI green}],
\end{equation}
\begin{equation}
\mathrm{apply\text{-}rate} = \frac{\#\{\text{patches that } \texttt{git apply}\}}{N_{\mathrm{att}}},
\qquad
\mathrm{coverage} = \frac{N_{\mathrm{att}}}{N}.
\end{equation}
This makes the honesty constraint explicit: a selective system can appear stronger on attempted cases partly by
declining the hard ones, so repair quality should never be read without its corresponding coverage term.

\subsection{Classification Metrics}

For classification studies, PatchOps reports per-class precision and recall through the standard \(F_1\) family.
For class \(c\),
\begin{equation}
F_{1,c} = \frac{2P_cR_c}{P_c + R_c},
\qquad
\mathrm{macro\text{-}}F_1 = \frac{1}{C}\sum_{c=1}^{C}F_{1,c},
\end{equation}
where \(C\) is the number of classes. Macro-\(F_1\) is useful here because weak categories should remain
visible instead of being dominated by larger easier classes.

\subsection{Decision and Abstention Metrics}

PatchOps is evaluated as a decision system, not only as a patch generator. If \(\hat{a}_i\) is the predicted
action and \(a_i^\star\) is the expected action label, act/hold accuracy is
\begin{equation}
\mathrm{Acc}_{\mathrm{act/hold}} = \frac{1}{N}\sum_i \mathbf{1}[\hat{a}_i = a_i^\star].
\end{equation}
Abstention accuracy uses the project's restricted denominator:
\begin{equation}
\mathrm{Acc}_{\mathrm{abst}} =
\frac{\#\{\text{should-abstain} \wedge \text{did abstain}\}}{\#\{\text{should-abstain}\}}.
\end{equation}
This definition matters because abstention is judged only on the subset where abstaining was the correct
behavior; it is not averaged over all decisions.

\subsection{Selective Prediction and Calibration}

For a confidence score \(s_i\), threshold \(\tau\), and zero-one loss \(\ell_i \in \{0,1\}\), selective
prediction is summarized through coverage, selective risk, and area under the risk-coverage curve:
\begin{equation}
\mathrm{cov}(\tau) = \frac{1}{N}\sum_i \mathbf{1}[s_i \ge \tau],
\qquad
\mathrm{risk}(\tau) = \frac{\sum_i \ell_i \mathbf{1}[s_i \ge \tau]}
{\sum_i \mathbf{1}[s_i \ge \tau]}.
\end{equation}
\begin{equation}
\mathrm{AURC} = \frac{1}{N}\sum_{k=1}^{N}\left(\frac{1}{k}\sum_{j=1}^{k}\ell_{(j)}\right),
\end{equation}
where \(\ell_{(j)}\) are losses sorted by decreasing confidence. Lower AURC indicates a better confidence
ranking because the accepted prefix stays more accurate over a wider range of coverage values.

Calibration is reported through expected calibration error. For \(M\) confidence bins \(B_m\),
\begin{equation}
\mathrm{ECE} = \sum_{m=1}^{M}\frac{|B_m|}{N}\left|\mathrm{acc}(B_m) - \mathrm{conf}(B_m)\right|.
\end{equation}
ECE measures whether stated confidence matches empirical correctness, which is why it is central to the
abstention story in Phase 11 and the learned decision gate in Phase 15.

\subsection{Plausibility and Ranking Quality}

Patch plausibility is evaluated through ROC AUC, which can be read as the probability that a random on-target
patch receives a higher score than a random off-target patch:
\begin{equation}
\mathrm{AUC} =
\frac{1}{|\mathcal{P}||\mathcal{N}|}
\sum_{p \in \mathcal{P}}\sum_{n \in \mathcal{N}} \mathbf{1}[s_p > s_n].
\end{equation}
This matches the role of the plausibility gate as a ranking mechanism rather than as a proof of semantic
correctness.

\subsection{Confidence Intervals}

Because several PatchOps studies operate at small \(n\) and sometimes near 0 or 1, Wilson score intervals are
preferable to normal approximations. For an empirical proportion \(\hat{p}\) and \(z = 1.96\),
\begin{equation}
\frac{\hat p+\frac{z^2}{2n}\;\pm\;z\sqrt{\frac{\hat p(1-\hat p)}{n}+\frac{z^2}{4n^2}}}{1+\frac{z^2}{n}}.
\end{equation}
This is the interval family behind the thesis wording that Phase 15 matches-or-beats the threshold baseline
rather than claiming a clean statistical separation at \(n=76\).

\section{Ground-Truth Design for Safety and Flakiness}

Safety and flakiness require more than a public benchmark label. The Safety Set therefore assigns structured
labels to each case, including expected action, risk class, allowed files, forbidden files, category, and
whether the case is flaky or transient. This allows PatchOps to be evaluated on the action it should have taken,
not merely on whether it happened to emit a patch.

This is methodologically important because the main PatchOps claim is asymmetric. On some failures, a patch is
the correct action. On others, not patching is the correct action. A benchmark that cannot represent that
distinction cannot evaluate the thesis contribution properly.

\section{Localization Methodology as a Separate Study}

The project treats localization as an explicit experimental subproblem. Rather than assuming that retrieved
context is always adequate, PatchOps compares log-path ranking, category reranking, LLM-guided candidates,
hybrid ranking, full-repository BM25, semantic embedding reranking, and later agentic exploration. All of these
studies reuse the same benchmark family and matching discipline, which makes it possible to interpret later
findings as a single evolving localization study rather than disconnected experiments.

This explicit localization design is one of the strongest methodological choices in the thesis. It enables later
claims such as ``localization is necessary but insufficient'' to be supported by measured transitions rather than
by informal intuition.

\section{Classifier Methodology as a Reproducibility Study}

The learned classifier phase is methodologically scoped as a reproducibility and cost study rather than as a raw
repair-improvement claim. The classifier is trained on compressed evidence features derived from the same
preprocessing pipeline the live engine uses. It is evaluated on held-out data under both exact and lenient
labeling rules and is compared with rule-based and LLM-based classification behavior.

This framing matters because it keeps the classifier's role proportional. It is not presented as the central
intelligence of PatchOps, but as a controlled and measurable way to make one stage of the pipeline open,
cheaper, and more reproducible.

\section{Selective Prediction and Calibration Methodology}

The selective prediction phase extends the evaluation from action accuracy to confidence quality. Here, the
question is whether the system's confidence signals can support principled abstention. The methodology compares
learned-classifier confidence against engine decision confidence on held-out data and evaluates whether acting on
only the more confident subset improves accuracy. Calibration methods such as isotonic regression are used to
measure whether these confidence values can be made more reliable without changing their ranking behavior.

This is important because it turns the thesis claim about abstention into something operationally tunable. A
system can then be evaluated not only at one fixed threshold, but along a family of possible operating points.

\section{Learned Decision-Gate Methodology}

The Phase 15 decision-gate study extends selective prediction from confidence analysis into the actual
act-versus-hold policy. The gate is trained and evaluated on the expanded 76-case Safety Set using
cross-validation. Its inputs are pre-generation signals that the engine already computes, including calibrated
classifier confidence, category risk, flaky status, semantic ambiguity, log-named-file count, and evidence
length. The comparison baseline is the earlier hand-tuned threshold policy.

This methodology is intentionally scoped as a decision-quality evaluation, not a repair-success evaluation. The
reported metrics are act/hold accuracy, abstention accuracy, over-act count, calibration error for
\(P(\mathrm{act})\), and risk-coverage behavior. The study also reports the low-risk lane separately, because
that is the branch where the deployed engine actually uses the learned gate.

Two caveats are built into the methodology. First, the expanded Safety Set is synthetic-heavy, so the result is
an indicative labeled-set result rather than a field measurement. Second, at \(n=76\), confidence intervals
overlap with the hand-tuned policy. The correct claim is therefore that the learned gate matches-or-beats the
threshold policy while removing the magic constants, not that it is statistically separated.

\section{Agentic Localization Methodology}

The agentic localization phase adds a bounded read-only navigation loop over a repository clone. This phase is
important methodologically because it tests whether the localization bottleneck is due merely to weak retrieval
or whether even tool-using exploration struggles to recover the right files. The experiment keeps the slice,
matching rules, and baseline comparisons aligned with earlier localization phases so that any gain or non-gain
is interpretable against the same benchmark context.

\section{Threats to Validity Built Into the Methodology}

Several threats to validity are acknowledged directly in the methodology rather than being postponed to the
results chapter alone.

First, some cells operate on modest sample sizes. The methodology therefore treats counts and coverage as part
of the result, not as footnotes.

Second, validation support is partial. Pass@1 style results are tied to the supported subset of workflows, and
that support boundary is disclosed rather than hidden.

Third, synthetic safety cases may be cleaner than reality. They are therefore disclosed explicitly and used for
the parts of the claim that require ground truth rather than for general repair claims.

Fourth, public benchmarks may be contaminated by model pretraining. PatchOps addresses this not by pretending
contamination is absent, but by scoping its strongest claims toward safety, abstention, and decision quality
instead of headline repair leadership.

\section{Methodological Honesty Rules}

Two practical rules govern the write-up of the later results.

First, when a figure or summary table disagrees with a per-phase findings document, the per-phase
\texttt{findings.md} takes precedence. This avoids circular drift from re-summarization.

Second, metrics must be read with their relevant context. For example, patch-apply rate must be interpreted
alongside coverage, because a selective system can improve apparent quality partly by refusing hard cases. The
project records these caveats as part of the methodology itself rather than treating them as optional
discussion.

\section{Chapter Conclusion}

This chapter described the evaluation methodology used to study PatchOps across diagnosis, localization, repair,
safety, abstention, and cost. The methodology combines a public CI repair benchmark, a labeled safety benchmark,
a bounded local validation harness, a modular ablation ladder, and a project-wide pre-registration discipline.
It also treats localization, classification, selective prediction, the learned decision gate, and agentic
navigation as explicitly measured subproblems rather than informal side effects of one large repair score.

With the datasets, baselines, metrics, and honesty rules now fixed, the next chapter can present the empirical
results themselves. Those results will be organized around the thesis's four headline findings, using the
canonical per-phase findings files as the source of truth.
